\section*{Chapitre \ref{ch:rs:stylised-facts-of-returns} --- Faits stylisés des rendements}

\begin{solution}{exo:rs:stylised-facts-of-returns:1}
$-17.4~\%/1.08~\% = -16.2$ écarts-types~; $\Phi(-16.2) \approx 6 \times 10^{-59}$, soit environ $10^{-54}$ jours de ce type sur 26\,296~: jamais, même sur un nombre quelconque
de durées de vie d'univers. Cela ne s'est produit qu'une fois.
\end{solution}

\begin{solution}{exo:rs:stylised-facts-of-returns:2}
Pour des termes indépendants, le kurtosis en excès se divise par le nombre de termes~: $3 + 17/5 = 6.4$, $3 + 17/21 = 3.8$, $3 + 17/63 = 3.3$.
Les valeurs observées sur le siècle, 11.5, 9.5 et 7.3, sont bien plus élevées~: les grands mouvements se regroupent dans le temps, de sorte qu'ils ne se compensent pas.
\end{solution}

\begin{solution}{exo:rs:stylised-facts-of-returns:3}
$\P(X > x) \propto x^{-\alpha}$, donc doubler le seuil divise la probabilité par $2^\alpha$~: avec $\alpha = 3$, un jour sur 8\,000 baisse de
plus de 8~\%~; avec $\alpha = 4$, un jour sur 16\,000.
\end{solution}

\begin{solution}{exo:rs:stylised-facts-of-returns:4}
Pour une loi de Student-$t$ à $\nu$ degrés de liberté, $\E[T^2] = \nu/(\nu - 2)$ et $\E[T^4] = 3\nu^2/((\nu - 2)(\nu - 4))$, donc le kurtosis est
$3(\nu - 2)/(\nu - 4) = 3 + 6/(\nu - 4)$, inchangé par un changement d'échelle. Un kurtosis de 19 exige $6/(\nu - 4) = 16$, d'où $\nu = 4.375$~: à peine plus de
quatre degrés de liberté.
\end{solution}

\begin{solution}{exo:rs:stylised-facts-of-returns:5}
La variance d'un portefeuille pondéré également de 10 actifs, de volatilité $\sigma$ et de corrélation deux à deux $\rho$, vaut $\sigma^2(1/10 +
0.9\rho)$~: $\sqrt{0.649} = 0.806\sigma$ les jours de baisse et $\sqrt{0.604} = 0.777\sigma$ les jours de hausse. La diversification supprime 19~\% de la
volatilité lorsque le marché baisse et 22~\% lorsqu'il monte.
\end{solution}

\begin{solution}{exo:rs:stylised-facts-of-returns:6}
Persistance $\alpha + \gamma/2 + \beta = 0.99$ (la moitié des chocs négatifs portent $\gamma$)~; un choc de variance se réduit de moitié en $\ln 0.5/\ln 0.99 = 69$
jours. L'autocorrélation des rendements absolus décroît comme la persistance élevée à la puissance du retard, exponentiellement~: faible après quelques
demi-vies, alors que la série réelle décroît comme une puissance du retard.
\end{solution}

\begin{solution}{exo:rs:stylised-facts-of-returns:7}
Avec $\alpha = 0.01$, $\gamma = 0.05$, $\beta = 0.965$, la persistance vaut exactement un (GARCH intégré)~: les autocorrélations de $|m|$
aux retards 20 et 100 atteignent 0.37 et 0.26, contre 0.21 et $-0.02$ pour la spécification par défaut. Mais la variance n'a plus de niveau vers lequel revenir
(l'ordonnée à l'origine $\omega$ est nulle)~: sur les dix années, la volatilité conditionnelle oscille entre 1.7~\% et 23.6~\% par an et la
volatilité réalisée vaut 7.7~\% au lieu des 16~\% visés. La mémoire longue exige un autre mécanisme (une composante de volatilité lente), et non une
racine unitaire.
\end{solution}

\begin{solution}{exo:rs:stylised-facts-of-returns:8}
Une couverture de queue rapporte lors des krachs de plusieurs jours et des phases de vente massive où les corrélations montent~; un simulateur calibré sur le kurtosis quotidien et
le regroupement de volatilité rate les deux (celui du livre reproduit les queues quotidiennes, mais affiche un kurtosis mensuel de 3.8 contre 9.5 et aucune \omterm{def:rs:stylised-facts-of-returns:exceedance}{asymétrie des corrélations}). Ses
événements de queue ne durent qu'un jour, trop courts et trop diversifiables, si bien que la valeur de la couverture est mal mesurée~; et des fenêtres de dix ans sur une seule
trajectoire simulée ne contiennent que très peu d'événements de queue. Il faut valider sur des données comportant des krachs, et sur un simulateur qui reproduit les queues mensuelles.
\end{solution}

\begin{solution}{pb:rs:stylised-facts-of-returns:1}
\textbf{1.} 19.1 et 20.5. \textbf{2.} 2.90 et 2.83. \textbf{3.} 103, contre 0.015. \textbf{4.} 22.0 (une baisse de 17.4~\% avec un
écart-type antérieur de 0.79~\%). \textbf{5.} 8.7~: les chocs propres à un titre ont des queues de Student-$t$ à quatre degrés de liberté
mais seule une partie de sa variance provient du facteur de marché, qui présente du regroupement. \textbf{6.} Données réelles 0.30, 0.22, 0.13~; marché synthétique 0.23,
0.21, $-0.02$. \textbf{7.} Au retard 100~: le GARCH décroît exponentiellement (demi-vie de 69 jours), la série réelle comme une puissance. \textbf{8.} Oui~:
$-0.111$ contre la valeur réelle $-0.085$. \textbf{9.} 9.5 contre 3.8~: les krachs du siècle ont duré des semaines (1929--1933, 1987, 2008) et
les régimes de volatilité ont duré des années~; le simulateur n'a ni l'un ni l'autre. \textbf{10.} Oui~: 0.22 pour $|r|$ contre 0.11 pour $r^2$.
\textbf{11.} 0.61 et 0.56. \textbf{12.} 0.20 les jours de baisse et 0.25 les jours de hausse~: le simulateur les classe dans le mauvais ordre, ce qui
sous-estime le coût d'un krach pour un portefeuille diversifié. \textbf{13.} 44.5, 22~\% du total, contre un edge de 3.58.
\textbf{14.} Quatre~: le marché et les facteurs sectoriels les plus puissants. \textbf{15.} Mesure des prédicteurs, mécanique des backtests, modèles de risque
aux horizons quotidien et hebdomadaire, construction de portefeuille, coûts. \textbf{16.} Couverture de queue, drawdowns sur plusieurs mois, diversification dans
les phases de vente massive, tout ce qui dépend d'une mémoire au-delà de cinquante jours. \textbf{17.} Krachs de plusieurs jours avec corrélations croissantes (bêtas
baissiers assortis d'une dérive compensatrice, ou un niveau de volatilité à changements de régime). \textbf{18.} \emph{Résultat nommé~:} le simulateur reproduit le
kurtosis quotidien réel (20.5 contre 19.1), l'indice de queue gauche (2.83 contre 2.90) et le regroupement de volatilité à vingt jours (0.21 contre 0.22)~;
il manque la mémoire à cent jours ($-0.02$ contre 0.13), le kurtosis mensuel (3.8 contre 9.5) et l'\omterm{def:rs:stylised-facts-of-returns:exceedance}{asymétrie des corrélations}
(0.20 à la baisse contre 0.25 à la hausse, là où les données donnent 0.61 contre 0.56). \textbf{19.} Parce qu'un résultat n'est aussi réaliste que les faits
que ses données reproduisent~; la grille d'évaluation indique quels résultats se transposent aux marchés. \textbf{20.} Il s'agit d'un test que tout modèle de rendements, ou tout
simulateur, doit réussir avant que l'on se fie à ses conclusions.
\end{solution}

\begin{solution}{iq:rs:stylised-facts-of-returns:1}
Queues épaisses (indice de queue d'environ 3)~; pas d'autocorrélation linéaire des rendements mais autocorrélation durable des rendements absolus
(regroupement de volatilité, mémoire longue)~; l'effet de levier~; \omterm{def:rs:stylised-facts-of-returns:aggregational}{gaussianité par agrégation} lente~; \omterm{def:rs:stylised-facts-of-returns:gainloss}{asymétrie gains--pertes}~; corrélations
qui montent lors des phases de vente massive~; \omterm{def:rs:stylised-facts-of-returns:intraday}{saisonnalité intrajournalière}.
\iqlookfor{plusieurs faits, avec un chiffre ou un mécanisme pour chacun.}
\end{solution}

\begin{solution}{iq:rs:stylised-facts-of-returns:2}
Non~: absence de corrélation ne signifie pas indépendance. Vérifier les autocorrélations de $|r|$ et de $r^2$ (regroupement de volatilité), un test d'effets ARCH,
et vérifier si le kurtosis décroît avec l'agrégation aussi vite que ne le prévoit l'indépendance.
\iqlookfor{la distinction, et un test concret de dépendance non linéaire.}
\end{solution}

\begin{solution}{iq:rs:stylised-facts-of-returns:3}
Les corrélations sont plus élevées les jours de baisse que les jours de hausse (0.61 contre 0.56 parmi les secteurs américains à un écart-type), et
la volatilité augmente avec elles, de sorte que la volatilité d'un portefeuille diversifié est plus élevée précisément lorsque surviennent les pertes~; le facteur de marché commun
domine lors d'une vente massive.
\iqlookfor{\omterm{def:rs:stylised-facts-of-returns:exceedance}{asymétrie des corrélations}, et pas seulement la hausse de la volatilité.}
\end{solution}

\begin{solution}{iq:rs:stylised-facts-of-returns:4}
Noter le simulateur sur les \omterm{def:rs:stylised-facts-of-returns:fact}{faits stylisés} qui comptent pour la stratégie~: les queues, le regroupement de volatilité et sa mémoire, l'effet de levier, la structure de corrélation et son
asymétrie, la \omterm{def:rs:stylised-facts-of-returns:intraday}{saisonnalité intrajournalière} si elle négocie en cours de journée, face à des données réelles~; puis vérifier que le résultat de la stratégie ne repose pas sur
un fait que le simulateur omet (faire varier ce trait et voir si la réponse bouge).
\iqlookfor{une grille d'évaluation liée aux expositions de la stratégie.}
\end{solution}

\begin{solution}{iq:rs:stylised-facts-of-returns:5}
$3 + 17/21 = 3.8$ sous indépendance. La valeur observée, 9, signifie que les rendements quotidiens sont dépendants par leur ampleur~: les grands mouvements arrivent
en grappes qui durent des semaines, si bien que les rendements mensuels héritent de queues épaisses~; c'est un regroupement de volatilité, et non seulement des queues quotidiennes épaisses.
\iqlookfor{la règle $1/n$ pour le kurtosis en excès et l'inférence sur la dépendance.}
\end{solution}

\begin{solution}{iq:rs:stylised-facts-of-returns:6}
La corrélation négative entre les rendements et la volatilité ultérieure~: les baisses font davantage monter la volatilité que les hausses. Mesurer la corrélation de
$r_t$ avec $r_{t+k}^2$ (ou avec la volatilité réalisée) pour $k > 0$, contre $k < 0$~; ou ajuster un GARCH asymétrique et tester $\gamma$.
\iqlookfor{le sens du temps dans la mesure.}
\end{solution}
